
over @@M0@@, since these points are in bijective correspondence with the flags of type @@M1@@ (@@M2@@ terms) of the vector space @@M3@@ of dimension @@M4@@ over @@M5@@. On the other hand, these rational points also correspond to the germs of sections of @@M6@@ in a neighborhood of @@M7@@, so one can find such a section @@M8@@ over an open subset @@M9@@ of @@M10@@. Let @@M11@@ be the closure of @@M12@@ with its reduced induced structure, and let @@M13@@ be the restriction of @@M14@@ to @@M15@@. The morphism @@M16@@ is projective and, since it is an isomorphism over @@M17@@, it is birational.
Thus @@M18@@ (EGA IV 5.6.6.1). On the other hand, @@M19@@ is closed and contains @@M20@@, so @@M21@@ is surjective and @@M22@@ (EGA IV 5.4.1(ii)).

The hypotheses of case I are therefore satisfied for @@M23@@ and @@M24@@, and one obtains
\sgaStableTarget{sga6.e10.eq.1-3-7.occ2}@@M25@@
Taking account of the projection formula and (1.1.8), one concludes that:
\sgaStableTarget{sga6.e10.eq.1-3-8}@@M26@@

On the other hand, the canonical homomorphism
@@M27@@
is an isomorphism over the dense open subset @@M28@@, and the @@M29@@-modules @@M30@@ (@@M31@@) are concentrated on the nowhere dense closed subset @@M32@@. It follows that
\sgaStableTarget{sga6.e10.eq.1-3-9}\begin{equation}\tag{1.3.9}
f_*(\cl_\bullet(\underline{O}_{X'}))=\cl_\bullet(\underline{O}_X)+w,
